\subsection{Integral coordinates and separation of rational frames}
\label{moment:rigidity}
Let $u=\varrho$, $n=u+1$, and write the ordered frame as
$F=(e_0,\ldots,e_{u-1},z)$, with $f_u=z$. In rational coordinates the Hermitian
space is the orthogonal sum of $u$ hyperbolic planes and $X_0\perp Ez$,
where $X_0$ has identity Gram matrix and $\langle z,z\rangle=-N$.
The $e_i$ pair with their duals $e_i^*$. At a complex place $v$, let
$\Phi_v$ be the chosen embedding and put $s_{va}=g_v^{u-a}/\sqrt N$ for
$0\le a\le u$, where $1\le |g_v|\le b^{23}$ is the local scale.
Changing to the normalized real-place basis multiplies the matrix entry
in row $i$, column $j$ by $s_{vj}/s_{vi}$ on frame rows and columns.
On $X_0$ use scale one and on $e_a^*$ the inverse scale of $e_a$.
The normalized Gram operator $J$ exchanges each hyperbolic pair, is the
identity on $X_0$ and acts by $-1$ on $z$. In particular $\|J\|=1$.

\paragraph{Two families of algebraic integers.}
Let $B_c$ be the finite adelic class basis. In the standard lattice basis
one middle column $x$ of $X_0$ is replaced by $x/\bar\pi+z/N$,
where $\pi\bar\pi=N$; the basis is then moved by the class representative
acting isometrically on $X_0$. Each $B_c$ fixes the frame columns and
leaves the rows along $e_i$ unchanged. Its integral Gram matrix
$J_{\rm lat}=B_c^*J_{\rm rat}B_c$ is independent of the class:
the altered column has squared Hermitian length $1/N-1/N=0$
and pairing $-1$ with $z$. For every rational term $g$ with
$K=B_c^{-1}gB_{c'}$ integral at the finite places, the two families
\[
 Z_{ja}(g)=\langle f_j,g f_a\rangle,\qquad
 A_{ia}(g)=(g f_a)_{e_i}\quad(0\le i<u)
\]
belong to $\mathcal O_E$. Indeed $Z_{ja}(g)$ is the $(f_j,f_a)$ entry
of $J_{\rm lat}K$ and $A_{ia}(g)$ is the corresponding row and
column of $K$. These equalities also hold at places dividing $N$.

\begin{lemma}\label{moment:fixed-frame}
Put $\rho=b^{37n}R$, $R\ge1$. In the normalized bases let $y_v$ be
points with $\|y_v^{-1}f_j\|\le\rho$ and Bessel length at most
$b^{6n+2}R$. Suppose that for a rational term $g$ as above, writing
$g_v$ for its normalized matrix,
\[
 \delta_{va}:=\|y_v^{-1}(g_v-1)f_a\|<(N\rho)^{-1}
 \qquad\text{for all }v,a.
\]
Then $gF=F$.
\end{lemma}
\begin{proof}
With the displayed change of basis,
\[
 |\Phi_v(Z_{ja}(g)-Z_{ja}(1))|
 \le |s_{vj}s_{va}|^{-1}\,\|y_v^{-1}f_j\|\,\delta_{va}
 \le N\rho\,\delta_{va}.
\]
This is Cauchy--Schwarz applied to the normalized Gram operator after
using invariance of the form under $y_v$. An algebraic integer of
modulus strictly below one at each chosen embedding is zero, by its
absolute norm. Hence every pairing agrees with its value for the identity.
Pairing with the $e_i$ and $z$ gives the expected dual and $z$
coordinates of $g f_a$. Unitarity gives the same length as $f_a$.
After those coordinates have been fixed, the difference can have no
$X_0$ component: the remaining contribution to the length is the sum
of its squared absolute coordinates at a complex embedding. Consequently
$(g-1)f_a=\sum_{i<u}d_i e_i$ with $d_i=A_{ia}(g)-A_{ia}(1)\in\mathcal O_E$.

Fix $a$ and descend on $i$. After $d_j=0$ for $j>i$, the $i$th row of
any Bessel unipotent $w$ applied to $(g_v-1)f_a$ equals
$\Phi_v(d_i)s_{va}/s_{vi}$, by unitriangularity on the isotropic flag.
The Bessel subgroup splits as a unipotent times a compact middle
isometry acting isometrically in the normalized Euclidean metric.
Thus the infimum defining the length permits a unipotent $w$ with
$\|w y_v\|\le b^{6n+2}R$ (the additive one in the length absorbs
the approximation to its infimum). It follows that
\[
 |\Phi_v(d_i)|\le b^{23u}\|w y_v\|\delta_{va}
 \le b^{23u+6n+2}R\,\delta_{va}\le\rho\,\delta_{va}<1 .
\]
The norm argument gives $d_i=0$ and closes the induction.
\end{proof}

\paragraph{Support and orbit weights.}
In rational frame coordinates take a reduced cutoff, supported where
the entries above the diagonal in the Gram--Schmidt coefficient matrix,
before archimedean rescaling, belong to a residue domain for
$M\mathcal O_E$ of sup radius $b^{26}$. At indices $a<j$ rescaling
multiplies the entry by $s_{vj}/s_{va}$ of modulus at most one.
The successive orthogonal residuals of $y_v^{-1}F$ have length at
most the Bessel length: the residual norms are unchanged by the
unitriangular action on the flag, and each is bounded by the norm of
the corresponding column after such a move. In normalized coordinates
$\|(wy_v)^{-1}\|=\|wy_v\|$ by $J$-unitarity. Thus each frame vector
has norm at most $n b^{26}\ell(y_v)\le b^{27}\ell(y_v)$.
On the length sublevel in \S\ref{B:cutoff} this is at most $\rho$.
For two contributing terms $p,q$ with the same starting class,
apply Lemma~\ref{moment:fixed-frame} to $p^{-1}q$ at $p^{-1}y$.
The relative term is integral between the two ending class bases since
the original finite matrices are invertible integral. Distinct rational
frames therefore have distance at least $(N\rho)^{-1}$ in the maximum
of the column norms over the embeddings.

If the frames and reductions agree, $p^{-1}q$ fixes $F$ and has
trivial finite reduction in the moved bases. Its component on
$X^\perp/X$ for $X=\langle e_i\rangle$ fixes $z$, hence acts on
$X_0$ by an isometry. The fitted level restricts on this definite
factor to the source level: the middle component on the quotient of
a principal congruence element lies in this level. It follows that
the ending representatives define the same source double coset.
Freeness at that level now makes the rational definite component
the identity. Thus the relative term belongs to the congruence
unipotents fixing the frame. Fixing one term in a fibre, every other
term gives a distinct such unipotent by taking the relative element.
For a cutoff whose translates by the class unipotent group sum to one,
nonnegativity bounds the weight of this fibre by one.

\paragraph{Construction of the cutoff on the required scales.}
We now construct the reduced cutoff just used, for the free level and
class bases of the fitted model. Write $\Gamma_c$ for the rational
Bessel unipotents in the principal congruence subgroup at $c$, and
$\Gamma_c^0$ for those fixing the frame pointwise. We compute the
image of $\Gamma_c$ on $F$. It consists of all upper unitriangular
$n$-by-$n$ matrices with entries in $M\mathcal O_E$ above the diagonal.
In fact $(g-1)/M$ is integral in the moved basis for $g\in\Gamma_c$.
At the split level primes it suffices to impose the congruence on the
chosen member of each pair: unitarity recovers the inverse adjoint
congruence on the conjugate member, as the lattice Gram matrix is
unimodular. At every other place $M$ is a unit. The top-row coordinates
computed above therefore give the inclusion in this unitriangular group.

For surjectivity fix $i<j$ and $a\in\mathcal O_E$. There is a rational
functional $\lambda$ integral on the class lattice, with
$\lambda(f_l)=a\,\mathbf1_{l=j}$ and zero on the other standard
basis vectors outside $X_0$. Indeed the corresponding adelic row of
$aB_c^{-1}$ has the prescribed entries there. Integrality on the
lattice is an open condition on the remaining finite-adelic entries,
so additive strong approximation supplies rational entries on $X_0$.
Write $\lambda=\langle v,\cdot\rangle$. The class lattice is self-dual
(the changed two-column Gram block has determinant $-1$), hence $v$
belongs to it. Put $x=e_i$, so that $\langle x,v\rangle=0$.
Choose $\tau=-M\zeta_3\in\mathcal O_E$, so $\tau+\bar\tau=M$.
Writing $v^\flat=\langle v,\cdot\rangle$, the map
\[
  g=1+M\bigl(xv^\flat-vx^\flat-\tau\langle v,v\rangle x x^\flat\bigr)
\]
preserves the form. In expanding $g^*J_{\rm rat}g$ the squared
term is $M^2\langle v,v\rangle (x^\flat)^*x^\flat$ and the linear
terms cancel it. The formula is integral and congruent to one at the
moved level. It fixes $e_l$ for $l\le i$, adds only multiples of $e_i$
to later isotropic frame vectors, and acts trivially on the middle
quotient; the adjoint action on the dual flag is then forced. Hence
it belongs to $\Gamma_c$. On the frame its matrix is $1+aM E_{ij}$
because $x^\flat$ vanishes on the frame. These elementary matrices
generate the indicated group by successive elimination of the upper
diagonals. The kernel is $\Gamma_c^0$.

The additive lattice $M\mathcal O_E$ in $\mathbb C^f$ has a Borel
domain of representatives of sup radius at most $b^{26}$.
For the fields of Theorem~\ref{thm:p1-data}, the $2f$ elements
$\zeta_{q_1}^a,\sqrt{-d}\,\zeta_{q_1}^a$, $0\le a<f$, are a
rational basis of algebraic integers by coprimality of the two
conductors. Their Minkowski images are a real basis. Rounding the
coordinates in this basis gives, for the sublattice generated by
their $M$-multiples, covering radius at most
$Mf(1+\sqrt d)/2$. The larger lattice has no greater covering
radius. Since $M\le b^{22}$, $f\le q_1\le b$ and $d$ is fixed,
the claim holds eventually. Choose a nearest representative in
the sup norm, breaking ties lexicographically among the finitely
many minimizing representatives. The resulting domain is Borel
by discreteness and countability.

Let $T(y)$ be the upper unitriangular matrix expressing the frame in
its Gram--Schmidt residuals, conjugated back by the diagonal scales.
When $y$ is replaced by $\gamma y$, $\gamma\in\Gamma_c$, this matrix
is right multiplied by the inverse frame matrix of $\gamma$. This
follows since adding earlier columns to a column preserves its
orthogonal residual; diagonal conjugation restores the original
rational coefficients. Reduce the entries of $T(y)$ into the chosen
additive domain in order of increasing superdiagonal. Right
multiplying by $1+M a E_{ij}$ changes entry $(i,j)$ by $Ma$,
affecting only more distant superdiagonals in that column. Existence
and uniqueness of each additive representative now give precisely one
coset of $\Gamma_c^0$ that takes $y$ into the reduced set.

Finally $\Gamma_c^0$ is discrete at infinity. A compact archimedean
set allows only finitely many rational matrices satisfying the fixed
finite-place lattice conditions, since their entries have bounded
denominators and bounded conjugates. There is a Borel fundamental
set for its free action by left translations: cover the ambient Lie
group by countably many open sets each disjoint from its distinct
translates and discard the orbits covered by earlier sets. Multiply
the indicator of this fundamental set by the reduced-set indicator.
Its sum over $\Gamma_c$ is exactly one, since reduction selected
one coset and the first indicator selects one element in that coset.
The same conclusion holds with any nonnegative measurable
$\Gamma_c^0$ cutoff whose translates sum to one. Thus the support
and weight properties hold for the class cutoffs used in the
amplitudes, regardless of that last choice.

On the constructed scales, F4 therefore follows for the fitted free
level, integral class bases, archimedean scales and reduced-cutoff
definition above. The existence of the actual filling with these
data (F0) and its identification with the period map (F1) are
separate requirements. The finite and archimedean basis computations
and the reduced-cutoff construction correspond to
\lineSrc{New/Spanning/Unfolding/Terms}{618}{frame integrality},
\lineSrc{New/Spanning/Unfolding/Packing}{286}{frame separation}, and
\lineSrc{New/Spanning/Unfolding/ReducedPoints}{637}{the box partition}.
